\chapterpageanex{Anexos}
\chapterstar{ANEXO 1: ESTUDIO DE VIABILIDAD Y FACTIBILIDAD.}

% Cambiar nombre de "Cuadro" a "Tabla"
\renewcommand{\tablename}{Tabla}
% Reiniciar contador de tablas
\setcounter{table}{0}
% Formato: A.1, A.2, A.3...
\renewcommand{\thetable}{A.\arabic{table}}

En el siguiente apartado se presentan las tablas que detallan los recursos esenciales que conforman el estudio de viabilidad y factibilidad del proyecto tecnológico, incluyendo los recursos materiales, tecnológicos, humanos y financieros necesarios para el desarrollo e implementación del sistema de gestión de inventario para la Farmacia Selene, con el propósito de mostrar de manera organizada los elementos requeridos para su elaboración.



\begin{longtable}{|p{3.5cm}|p{5cm}|p{5.5cm}|}

\caption{Recursos Humanos} \label{tab:recursos_humanos} \\

\hline
\rowcolor{gray!20}
\textbf{Puesto} & \textbf{Perfil} & \textbf{Función} \\
\hline
\endfirsthead

\hline
\rowcolor{gray!20}
\textbf{Puesto} & \textbf{Perfil} & \textbf{Función} \\
\hline
\endhead

\hline
\multicolumn{3}{|r|}{\small Continúa en la siguiente página} \\
\hline
\endfoot

\hline
\endlastfoot

Desarrollador Backend & Conocimientos en programación con Python, desarrollo de APIs REST con FastAPI y manejo de bases de datos relacionales. & Diseñar e implementar la lógica de negocio del sistema, estructurar la base de datos y garantizar la correcta comunicación entre frontend, dispositivos IoT y servidor. \\
\hline

Desarrollador Frontend & Conocimientos en HTML, CSS, JavaScript y React, así como en el uso de bibliotecas de componentes como Material-UI. & Diseñar e implementar la interfaz gráfica del sistema, asegurando una experiencia de usuario clara, funcional y responsiva. \\
\hline

Desarrollador de Sistemas Embebidos & Conocimientos en programación en C++, manejo de microcontroladores y protocolos de comunicación como I²C y HTTP. & Programar el microcontrolador ESP32 para la lectura de sensores, procesamiento básico de datos y envío de información al backend. \\
\hline

Administrador de Base de Datos & Conocimientos en modelado de bases de datos, consultas SQL y administración de PostgreSQL. & Diseñar la estructura de la base de datos, validar la integridad de la información y supervisar el almacenamiento de datos del sistema. \\
\hline

Personal de Farmacia & Conocimiento operativo del inventario, ventas y manejo de medicamentos. & Utilizar el sistema para registrar ventas, consultar inventario y atender alertas generadas por el monitoreo ambiental. \\
\hline

\end{longtable}






\begin{longtable}{|p{3.5cm}|p{2.5cm}|p{8cm}|}

\caption{Recursos Materiales} \label{tab:recursos_materiales} \\

\hline
\rowcolor{gray!20}
\textbf{Componente} & \textbf{Cantidad} & \textbf{Especificaciones} \\
\hline
\endfirsthead

\hline
\rowcolor{gray!20}
\textbf{Componente} & \textbf{Cantidad} & \textbf{Especificaciones} \\
\hline
\endhead

\hline
\multicolumn{3}{|r|}{\small Continúa en la siguiente página} \\
\hline
\endfoot

\hline
\endlastfoot

Microcontrolador ESP32-C5 & 1 & Microcontrolador con conectividad Wi-Fi 6 (802.11ax) y Bluetooth Low Energy (BLE), utilizado para la lectura de sensores y comunicación con el backend. \\
\hline

Sensor AHT10 & 1 & Sensor digital de temperatura (±0.3 °C) y humedad relativa (±2\%), comunicación mediante protocolo I²C. \\
\hline

Lector de códigos de barras & 1 & Dispositivo de entrada con conexión USB o serial para captura automática de códigos de productos. \\
\hline

Cables Dupont & 1 juego & Cables macho-macho y macho-hembra utilizados para la conexión entre el ESP32 y el sensor AHT10. \\
\hline

Fuente de alimentación USB & 1 & Adaptador de corriente 5V para alimentación del microcontrolador ESP32. \\
\hline

Cable USB & 1 & Cable USB utilizado para programación y alimentación del microcontrolador. \\
\hline

Computadora de desarrollo & 1 & Equipo de cómputo con capacidad para ejecutar herramientas de desarrollo (IDE, base de datos y navegador web). \\
\hline

Router / Red Wi-Fi & 1 & Infraestructura de red local necesaria para la comunicación entre dispositivos IoT, backend y frontend. \\
\hline

\end{longtable}









\begin{longtable}{|p{3.5cm}|p{8cm}|p{2.5cm}|}

\caption{Recursos Tecnológicos} \label{tab:recursos_tecnologicos} \\

\hline
\rowcolor{gray!20}
\textbf{Recurso Tecnológico} & \textbf{Descripción} & \textbf{Tipo} \\
\hline
\endfirsthead

\hline
\rowcolor{gray!20}
\textbf{Recurso Tecnológico} & \textbf{Descripción} & \textbf{Tipo} \\
\hline
\endhead

\hline
\multicolumn{3}{|r|}{\small Continúa en la siguiente página} \\
\hline
\endfoot

\hline
\endlastfoot

Python & Lenguaje de programación utilizado para el desarrollo del backend y la lógica de negocio del sistema. & Software \\
\hline

FastAPI & Framework web empleado para la construcción de la API REST que gestiona la comunicación entre frontend, dispositivos IoT y base de datos. & Software \\
\hline

PostgreSQL & Sistema de gestión de base de datos relacional utilizado para almacenar información de usuarios, productos, ventas y registros ambientales. & Software \\
\hline

SQLAlchemy & Herramienta ORM utilizada para interactuar con la base de datos mediante estructuras del lenguaje Python. & Software \\
\hline

React & Biblioteca de JavaScript utilizada para la construcción de la interfaz de usuario del sistema. & Software \\
\hline

Material-UI (MUI) & Biblioteca de componentes visuales que permite mantener consistencia y diseño responsivo en la interfaz web. & Software \\
\hline

Vercel & Plataforma en la nube utilizada para el despliegue y alojamiento del frontend. & Servicio en la nube \\
\hline

Railway & Plataforma de infraestructura en la nube utilizada para el despliegue del backend y la base de datos. & Servicio en la nube \\
\hline

GitHub & Plataforma de control de versiones utilizada para el almacenamiento y gestión del código fuente del proyecto. & Servicio en la nube \\
\hline

Visual Studio Code & Entorno de desarrollo utilizado principalmente para la construcción del frontend con React y la gestión de archivos del proyecto. & Software \\
\hline

PyCharm & Entorno de desarrollo integrado utilizado para la programación, depuración y estructuración del backend desarrollado con Python y FastAPI. & Software \\
\hline

DataGrip & Herramienta de administración y visualización de bases de datos utilizada para consultar, validar y gestionar la información almacenada en PostgreSQL. & Software \\
\hline

\end{longtable}









\xdef\TotalGeneral{0}

\begin{longtable}{|p{4.2cm}|S|r|S|}
\caption{Recursos Financieros} \label{tab:recursos_financieros} \\

\hline
\rowcolor{gray!20}
\textbf{Componente} & 
\textbf{Precio Unitario (MXN)} & 
\textbf{Cantidad} & 
\textbf{Total (MXN)} \\
\hline
\endfirsthead

\hline
\rowcolor{gray!20}
\textbf{Componente} & 
\textbf{Precio Unitario (MXN)} & 
\textbf{Cantidad} & 
\textbf{Total (MXN)} \\
\hline
\endhead

\hline
\multicolumn{4}{|r|}{\small Continúa en la siguiente página} \\
\hline
\endfoot

\hline
\endlastfoot

\fila{ESP32-C5 Microcontrolador}{80}{1}
\fila{Sensor AHT10}{150}{1}
\fila{Lector de códigos de barras}{350}{1}
\fila{Cables Dupont}{50}{1}
\fila{Fuente de alimentación USB}{150}{1}
\fila{Cable USB}{50}{1}
\fila{Computadora de desarrollo}{10000}{1}
\fila{Router / Red Wi-Fi local}{800}{1}
\fila{Vercel (Hosting Frontend, Pro)}{360}{1}[1 mes]   % <-- uso del opcional
\fila{Railway (Hosting Backend + DB)}{360}{1}[1 mes]   % <-- uso del opcional

\multicolumn{3}{|r|}{\textbf{Total general estimado}} &
\num{\TotalGeneral} \\

\end{longtable}


\section*{Impacto Social del Proyecto}

El desarrollo e implementación del sistema de gestión de inventario para la Farmacia Selene genera un impacto social positivo tanto a nivel organizacional como comunitario. La digitalización de procesos administrativos y el monitoreo automatizado de condiciones ambientales contribuyen a mejorar la eficiencia operativa y la calidad del servicio ofrecido a los clientes.

En primer lugar, el sistema reduce errores humanos asociados al registro manual de ventas y control de inventario, lo que disminuye pérdidas económicas y mejora la disponibilidad de medicamentos. Esto impacta directamente en la comunidad, ya que garantiza un acceso más confiable y oportuno a productos farmacéuticos esenciales.

En segundo lugar, la incorporación de sensores de temperatura y humedad permite supervisar de manera continua las condiciones de almacenamiento de medicamentos sensibles. Este monitoreo ayuda a preservar la calidad y eficacia de los productos, protegiendo la salud de los pacientes y promoviendo buenas prácticas de conservación farmacéutica.

Desde el punto de vista laboral, el sistema no reemplaza al personal, sino que optimiza sus actividades, permitiéndoles concentrarse en tareas de atención al cliente y gestión administrativa de mayor valor. Además, fomenta la adopción de herramientas tecnológicas en pequeñas empresas, impulsando la transformación digital en el sector comercial local.

Finalmente, el proyecto demuestra que tecnologías modernas como aplicaciones web, bases de datos en la nube y dispositivos IoT pueden integrarse en negocios de pequeña escala de manera accesible y sostenible. Esto abre la posibilidad de replicar el modelo en otras farmacias o establecimientos similares, contribuyendo al fortalecimiento tecnológico del entorno empresarial y mejorando la calidad del servicio en la comunidad.